\exercisegroup

\begin{enumerate}
\def\labelenumi{\arabic{enumi})}
\item
  Let \(E_{0}=Q_{p}^{N}\) be the vector space over the p-adic field \(Q_{p}\) (GT, III, § 6, exerc. 23) which is the product of an enumerable infinity of factor each identical with \(Q_{p}\) . Let \(P\subset E_{0}\) be the set \(Z_{p}^{N}\) , and let E be the vector subspace of \(E_{0}\) generated by P. On the additive group P we consider the product compact topology of the topologies of the factors \(Z_{p}\) , and we denote by B the filter of neighbourhoods of 0 in P for this topology. Show that B is a fundamental system of neighbourhoods of 0 in E for a topology T compatible with the additive group structure of E, that satisfies \((\mathrm{EVT}_{\mathrm{I}}^{\prime})\) and \((\mathrm{EVT}_{\mathrm{III}}^{\prime})\) but not \((\mathrm{EVT}_{\mathrm{II}}^{\prime})\) (prove that the homothety \(x\mapsto x/p\) is not continuous in E).
\item
  Let K be a non-discrete topological division ring, \(K_{0}\) the division ring K with the discrete topology. The discrete topology on \(K_{0}\) is compatible with its additive group structure, and when we consider \(K_{0}\) as a vector space over K, it satisfies the axioms (EVT \(_{II}\) ) and (EVT \(_{III}\) ) but not (EVT \(_{I}\) ).
\item
  For every real number \(\alpha > 0\) , let \(G_{\alpha}\) be the topological group \(\mathbf{R} / \alpha \mathbf{Z}\) , and let \(G\) be the topological product group \(\prod_{\alpha} G_{\alpha}\) ( \(\alpha\) varying in the set of number \(>0\) ). For every \(x \in \mathbf{R}\) ,
\end{enumerate}

let \(t_{\alpha}(x)\) be the canonical image of \(x\) in \(G_{\alpha}\) ; the mapping \(\phi : x \to (t_{\alpha}(x))\) is a continuous injective homomorphism of \(\mathbf{R}\) in \(G\) . We consider on \(\mathbf{R}\) the topology that is the inverse image by \(\phi\) of that of \(G\) , and denote by \(E\) the topological group formed by \(\mathbf{R}\) with this topology. Show that when \(E\) is considered as a vector space over \(\mathbf{R}\) its topology satisfies \((\mathrm{EVT}_{I}^{\prime})\) and \((\mathrm{EVT}_{II}^{\prime})\) but not \((\mathrm{EVT}_{III}^{\prime})\) .

\begin{enumerate}
\def\labelenumi{\arabic{enumi})}
\setcounter{enumi}{3}
\item
  Let E be a vector space over a division ring K with a valuation; we suppose that E carries a metrisable topology compatible with its additive group structure. Suppose further that this topology satisfies axioms (EVT \(_{I}\) ) and (EVT \(_{II}\) ); show that if one of the two metrisable groups K, E is complete then the topology also satisfies (EVT \(_{III}\) ) and is, in consequence, compatible with the vector space structure of E (cf.~GT, IX, § 5, exerc. 23).
\item
  Let \(\mathbf{K}\) be a non-discrete valued field, and \(S\) be an arbitrary infinite set.
\end{enumerate}

\begin{enumerate}
\def\labelenumi{\alph{enumi})}
\item
  Let \(\mathbf{D} = (a_n)\) be an enumerable infinity of elements of \(\mathbf{S}\) . For every \(\lambda \in \mathbf{K}\) such that \(|\lambda| \leqslant 1\) , let \(u_{\lambda}\) be the element of the normed space \(\mathcal{B}_{\mathbf{K}}(\mathbf{S})\) , (I, p.~4, Example) of bounded mappings of \(\mathbf{S}\) in \(\mathbf{K}\) , such that \(u_{\lambda}(a_n) = \lambda^n\) for all \(n \in \mathbf{N}\) and \(u_{\lambda}(b) = 0\) for \(b \notin \mathbf{D}\) . Show that the family \((u_{\lambda})\) is algebraically independent.
\item
  Deduce that every basis of the vector space \(\mathcal{B}_{\mathbf{K}}(\mathbf{S})\) has the same cardinality as \(\mathrm{K}^{\mathrm{s}}\) (using \(a\) ) show that the cardinal of every basis of \(\mathcal{B}_{\mathbf{K}}(\mathbf{S})\) is at least equal to Card (K); note on the other hand that \(\operatorname{Card}(\mathcal{B}_{\mathbf{K}}(\mathbf{S})) = \operatorname{Card}(\mathbf{K}^{\mathrm{s}})\) and use A, II, § 2, exerc. 22).
\item
  Show in the same way that every basis of the vector space \(\ell_{\mathbf{K}}^{1}(\mathbf{S})\) has the cardinality of \((\mathbf{K}\times \mathbf{S})^{\mathbf{N}}\) .
\end{enumerate}

\begin{enumerate}
\def\labelenumi{\arabic{enumi})}
\setcounter{enumi}{5}
\tightlist
\item
  Let \(\mathbf{K}\) be a non-discrete valued division ring. Show that, for the space \(\ell_{\mathbf{K}}^{1}(\mathbf{N})\) of absolutely summable sequences \(x = (\xi_n)\) of elements of \(\mathbf{K}\) , the norms \(\| x \|_1 = \sum_{n=0}^{\infty} |\xi_n|\) and \(\| x \| = \sup_n |\xi_n|\) are not equivalent (cf.~GT, IX, § 3.3, prop. 7); show that \(\ell_{\mathbf{K}}^{1}(\mathbf{N})\) with the norm \(\| x \|\) is never complete even if \(\mathbf{K}\) is complete; what is its closure in \(\mathcal{B}_{\mathbf{K}}(\mathbf{N})\) ?
\end{enumerate}

¶ 7) * Let A be a ring with a discrete valuation, v the normed valuation of the division ring of fractions K of A; take the absolute value \(a^{v}\) on K, where 0 \textless{} a \textless{} 1. Let E be a normed vector space over K, for which the norm satisfies the ultrametric inequality

\[
\| x + y \| \leqslant \sup (\| x \|, \| y \|).
\]

\begin{enumerate}
\def\labelenumi{\alph{enumi})}
\item
  Denote by M the set of those \(x \in E\) for which \(\| x \| \leqslant 1\) , and by \(\pi\) a uniformizer of A; M is an A-module, and \(M / \pi M\) a vector space on the residual division ring \(k = A / \pi A\) of A. Let \((e_{\lambda})_{\lambda \in L}\) be a family of elements of M such that the images of \(e_{\lambda}\) in \(M / \pi M\) form a basis of this vector \(k\) -space. Show that \((e_{\lambda})\) is an independent family in E and that the vector subspace F of E generated by \((e_{\lambda})\) is dense in E.
\item
  If we put \(\| x\| _1 = \sup_{\lambda}|\xi_\lambda |,\) for every \(x = \sum_{\lambda}\xi_{\lambda}e_{\lambda}\) in F, show that on F the norms \(\| x\|\) and \(\| x\| _1\) are equivalent.
\item
  Let \(\mathbf{K}\) be complete. Deduce from \(a)\) and \(b)\) that, if \(\mathbf{L}\) is finite, the completion \(\hat{\mathbf{E}}\) of \(\mathbf{E}\) is isomorphic to \(\mathbf{K}^{\mathrm{L}}\) ; if \(\mathbf{L}\) is infinite \(\hat{\mathbf{E}}\) is isomorphic to the subspace \(\mathcal{C}_{\mathbf{K}}^{0}(\mathbf{L})\) of \(\mathcal{B}_{\mathbf{K}}(\mathbf{L})\) formed of the families \((\xi_{\lambda})\) such that \(\lim \xi_{\lambda} = 0\) for the filter of complements of finite subsets of \(\mathbf{L}\) .
\item
  We suppose K and E complete; let G be a second normed complete space over K whose norm satisfies the ultrametric inequality. Show that on replacing (if necessary) the norm of \(\mathcal{L}(\mathrm{E};\mathrm{G})\) (GT, X, § 3.2) by an equivalent one, then \(\mathcal{L}(\mathrm{E};\mathrm{G})\) is isometric to the vector space of families \((y_{\lambda})_{\lambda \in L}\) of elements of G such that \(\sup_{\lambda \in L}\| y_{\lambda}\| < +\infty\) , carrying the norm \(\sup_{\lambda \in L}\| y_{\lambda}\|\) (which is also an ultrametric norm).
\end{enumerate}

\begin{enumerate}
\def\labelenumi{\arabic{enumi})}
\setcounter{enumi}{7}
\item
  Let E be a topological vector space over a non-discrete topological division ring K. In order that there should exist a neighbourhood of the point \((0,0)\) in \(K \times E\) such that the mapping \((\lambda, x) \mapsto \lambda x\) should be uniformly continuous in this neighbourhood, it is necessary and sufficient that there exist a neighbourhood \(V_{0}\) of 0 in E such that the sets \(\lambda V_{0}\) form a fundamental system of neighbourhoods of 0 in E, where \(\lambda\) varies in the set of elements \(\neq 0\) of K. When K is a division ring with a non-discrete valuation and E is Hausdorff, show that the uniform structure of E is then metrisable.
\item
  Generalize prop. 5 of I, p.~8, to the case where the spaces \(\mathrm{E}_i\) ( \(\leqslant i \leqslant n\) ) and \(\mathrm{F}\) are topological vector spaces over an arbitrary non-discrete topological field.
\item
  Let E be a complete Hausdorff topological vector space over a non-discrete valued division ring K. Denote by F a vector subspace of E, and by T the topology on F induced by the topology \(T'\) of E; let B be a fundamental system of closed, balanced neighbourhoods of 0 for the topology T. Let \(F_{0}\) be the vector subspace of E, generated by the closures \(\overline{V}\) in E (relative to \(T'\) ) of the sets \(V \in B\) ; the sets \(\overline{V}\) form a fundamental system of neighbourhoods of 0 for a topology \(T_{0}\) on \(F_{0}\) , compatible with the vector space structure of \(F_{0}\) ; for this topology, \(F_{0}\) is complete, and the topology induced by \(T_{0}\) on F is identical with T.
\item
  In a topological vector space E over a non-discrete topological division ring K there exists a fundamental system B of closed neighbourhoods of 0, satisfying the conditions (EV \(_{II}\) ) and (EV \(_{III}\) ) as well as the two following :
\end{enumerate}

\((\mathrm{EV}_{\mathrm{Ia}})\) For each \(\mathbf{V} \in \mathfrak{B}\) , there exists \(\mathbf{W} \in \mathfrak{B}\) and a neighbourhood \(\mathbf{U}\) of 0 in \(\mathbf{K}\) such that \(\mathrm{UW} \subset \mathrm{V}\) . \((\mathrm{EV}_{\mathrm{Ib}})\) For every \(x \in \mathbf{E}\) and every \(\mathbf{V} \in \mathfrak{B}\) , there exists \(\lambda \neq 0\) in \(\mathbf{K}\) such that \(\lambda x \in \mathbf{V}\) . Conversely, let \(\mathbf{E}\) be a vector space over \(\mathbf{K}\) and let \(\mathfrak{B}\) be a filter base on \(\mathbf{E}\) satisfying the conditions \((\mathrm{EV}_{\mathrm{Ia}}), (\mathrm{EV}_{\mathrm{Ib}}), (\mathrm{EV}_{\mathrm{II}})\) and \((\mathrm{EV}_{\mathrm{III}})\) . Show that there is a topology on \(\mathbf{E}\) (and one such only), that is compatible with the vector space structure of \(\mathbf{E}\) , and for which \(\mathfrak{B}\) is a fundamental system of neighbourhoods of 0.

\begin{enumerate}
\def\labelenumi{\arabic{enumi})}
\setcounter{enumi}{11}
\item
  Let K be a discrete field, E the division ring of fractions of the ring of formal series A = K{[}{[}X, Y{]}{]} in two indeterminate variables on K (A, IV, p.~36). For every \(n \geqslant 0\) , let \(V_{n} \subset A\) be the set of formal series of total degree at least equal to n.~Show that in E, the sets \(V_{n}\) form a fundamental system of neighbourhoods of 0, for a topology compatible with the vector space structure of E (over K), for which E is metrisable and complete; if further K is a finite field, then E is locally compact. Show that the K-bilinear mapping \((u, v) \mapsto uv\) of E × E in E is continuous at the point \((0, 0)\) but that there exists \(u_{0} \in E\) such that \(v \mapsto u_{0}v\) is not continuous in E (for example \(u_{0} = 1/X\) ).
\item
  Let E be a vector space of infinite dimension over R, and let T be the family of all absorbent and balanced sets of E. Show that T does not satisfy axiom \((\mathrm{EV}_{\mathrm{III}})\) (in other words is not a fundamental system of neighbourhoods of 0 for a topology compatible with the additive group structure of E). For this, consider an infinite independent family \((e_{n})_{n\geqslant1}\) in E; for every integer \(n\geqslant1\) , let \(A_{n}\) be the set of points \(\sum_{i=1}^{n}t_{i}e_{i}\) such that \(|t_{i}|\leqslant1/n\) for \(1\leqslant i\leqslant n\) ; let A be the union of the \(A_{n}\) , and V be a subspace complementary to the subspace of E generated by the \(e_{n}\) , and write C for the set \(A+V\) ; show that there exists no set \(M\in T\) such that \(M+M\subset C\) .
\end{enumerate}

¶ 14) Let K be a Hausdorff topological division ring, \((E_{t})_{t\in I}\) an infinite family of Hausdorff topological vector spaces on K, none of which is the single point 0. We consider on \(F = \prod_{t\in I} E_t\) the topology \(\mathcal{T}\) , compatible with the additive group structure of \(F\) , for which a fundamental system of neighbourhoods of 0 is formed by the products \(\prod_{\iota \in I} V_{\iota}\) , where, for each \(\iota \in I\) , the set \(V_{t}\) is a neighbourhood of 0 in E (this topology is strictly finer than the product topology; cf.~GT, III, § 2, exerc. 23). We denote by \(T_{0}\) the topology induced by T on the subspace \(E = \bigoplus_{i \in I} E_{i}\) of F; E is closed in F for the topology T, and if each of the \(E_{t}\) is complete, then F is complete for the topology T, therefore E is complete for the topology \(T_{0}\) (GT, III, § 3, exerc. 10).

\begin{enumerate}
\def\labelenumi{\alph{enumi})}
\item
  Show that if there exists in K a neighbourhood of 0 bounded on the right (GT, III, § 6, exerc. 12) (in particular if K is a division ring with a valuation), the topology \(\mathcal{T}_0\) is compatible with the vector space structure of E. If, further, K is not discrete, then E is not a Baire space for any topology that is finer than \(\mathcal{T}_0\) and compatible with the vector space structure of E. b) Moreover, if there does not exist in K any neighbourhood of 0 bounded on the right (see c)) give an example of a family \((\mathrm{E}_t)\) such that the topology \(\mathcal{T}_0\) is not compatible with the vector space structure of E.
\item
  Let \(\mathbf{A} = \mathbf{R}[\mathbf{X}]\) be the ring of polynomials in one variable on \(\mathbf{R}\) . For every sequence \(s = (\varepsilon_n)_{n\geqslant 0}\) of real numbers \(>0\) , denote by \(\mathrm{V}_s\) the set of polynomials \(\sum_k a_k\mathrm{X}^k\in \mathrm{A}\) such that \(|a_k| < \varepsilon_k\) for all k. Let T be the set of the \(V_{s}\) where s varies in the set of sequences of numbers \textgreater0. Show that T is a fundamental system of symmetric neighbourhoods of 0 for a topology compatible with the ring structure of A. Let \(\mathbf{K} = \mathbf{R}(\mathbf{X})\) be the division ring of fractions of A; denote by S the family of subsets of K of the form \(\mathrm{U}(1 + \mathrm{U})^{-1}\) , where U varies in the set of the \(V_{s}\) not containing 1; show that S is a fundamental system of neighbourhoods of 0 for a topology compatible with the division ring structure of K, and that there does not exist in K any neighbourhood of 0 that is bounded.
\end{enumerate}

\(d\) ) For every Hausdorff topological division ring \(\mathbf{K}\) , show that there exists a set I such that on \(\mathbf{F} = \mathbf{K}^{\mathrm{I}}\) , the topology \(\mathcal{T}\) , defined above, is not compatible with the vector space structure of F.
% semantic-fragments: legacy-input-seam
\relax{}
